% Qatra — Getting the First Committee
%
% Compile with XeLaTeX (Arabic in the title):
%   xelatex -output-directory=build docs/qatra-pilot-guide.tex
%
% A companion to the dossier, and deliberately a different document. The
% dossier is for the person being presented to. This is for the person walking
% into a room to ask for something.
\documentclass[11pt,a4paper]{article}

\usepackage{fontspec}
\usepackage[margin=2.4cm]{geometry}
\usepackage{xcolor}
\usepackage{booktabs}
\usepackage{tabularx}
\usepackage{enumitem}
\usepackage{titlesec}
\usepackage{fancyhdr}
\usepackage[hidelinks]{hyperref}
\usepackage{parskip}

\definecolor{qatraRed}{HTML}{E5484D}
\definecolor{qatraNavy}{HTML}{0B2432}
\definecolor{qatraTeal}{HTML}{0E8BA8}
\definecolor{qatraGrey}{HTML}{5A6B75}
\definecolor{qatraGreen}{HTML}{12B76A}
\definecolor{qatraAmber}{HTML}{F5871F}
\definecolor{sayBg}{HTML}{F4FAFB}
\definecolor{rule}{HTML}{DCE4E7}

\setmainfont{Segoe UI}
\newfontfamily\arabicfont[Script=Arabic]{Segoe UI}
\newcommand{\ar}[1]{{\arabicfont #1}}

\titleformat{\section}{\Large\bfseries\color{qatraNavy}}{\thesection.}{0.6em}{}
\titleformat{\subsection}{\large\bfseries\color{qatraTeal}}{\thesubsection}{0.6em}{}
\titlespacing*{\section}{0pt}{1.4em}{0.6em}

\pagestyle{fancy}
\fancyhf{}
\fancyhead[L]{\small\color{qatraGrey}Qatra — Getting the First Committee}
\fancyhead[R]{\small\color{qatraGrey}\thepage}
\renewcommand{\headrulewidth}{0.4pt}

\newcommand{\note}[1]{{\small\color{qatraGrey}\textit{#1}}}

% Something to actually say out loud. Built from a saved box rather than
% mdframed, which patches the same table internals bidi does.
\newsavebox{\saybox}
\newenvironment{say}
  {\par\medskip\noindent\begin{lrbox}{\saybox}%
   \begin{minipage}{\dimexpr\textwidth-2\fboxsep-2pt}\itshape\vspace{4pt}}
  {\vspace{4pt}\end{minipage}\end{lrbox}%
   \fcolorbox{rule}{sayBg}{\usebox{\saybox}}\par\medskip}

\begin{document}

\begin{titlepage}
\centering
\vspace*{3cm}
{\Huge\bfseries\color{qatraNavy} Qatra \ar{قطرة}}\\[0.5em]
{\Large\color{qatraRed} Getting the First Committee}\\[2em]
{\large\color{qatraGrey} One committee. Twenty donors. One afternoon.}\\[3em]
\note{A companion to the technical dossier, and a different kind of document.\\
The dossier is written for the person being presented to.\\
This one is written for the person walking into a room to ask for something.}\\[2em]
{\small\color{qatraGrey} 1 September 2026}
\end{titlepage}

\tableofcontents
\newpage

% ===========================================================================
\section{What the problem actually is}

Everything else is finished. The loop closes, the law is respected, the tests
pass, the site is live and installable. The one thing missing cannot be built:
\textbf{nobody uses it}.

A panel will ask. The answer today is \emph{nobody yet}, and that answer is
survivable exactly once — if it is followed immediately by what is being done
about it. This document is what is being done about it.

\subsection{Why twenty is the right number}

Not because twenty is impressive. Because twenty is the smallest number that
turns every architectural claim in the dossier into evidence.

With twenty real donors in one wilaya, a request posted at nine in the evening
reaches people who can actually answer it. Without them, the same request
reaches nobody and proves nothing. The application does not become true at a
thousand users. It becomes true at the first one who receives a notification
and goes.

\subsection{Why it is not as hard as it sounds}

You are not recruiting twenty strangers. You are asking one person who
\emph{already has a list} to point it at something. The Croissant-Rouge
committees, the university clubs and the scout groups have been organising
blood donation in Algeria for decades. None of them lacks donors. What they
lack is a way to reach the right ones on a Tuesday night.

% ===========================================================================
\section{Who to ask}

\begin{center}
\begin{tabularx}{\textwidth}{@{}rlX@{}}
\toprule
& \textbf{Who} & \textbf{Why them} \\
\midrule
1 & A wilaya committee of the \textbf{FADS} &
The F\'ed\'eration Alg\'erienne des Donneurs de Sang is a federation of
blood-donor associations with local committees. Its entire reason to exist is
donor lists --- which is the one thing this needs. The ANS names it first among
its association partners. \\
2 & A university blood-donation club &
The shortest distance if you are a student. The president already has a group
of donors on their phone, and nobody has to ask a head office. Fastest yes on
this list. \\
3 & A local \textbf{Croissant-Rouge} committee &
Carries the most weight in front of a panel, and the ANS lists the CRA as a
collection partner in writing. Slower: somebody must say yes on the
organisation's behalf. \\
4 & \textbf{Scouts Musulmans Alg\'eriens} &
Also named on the ANS partnership page, and historically active in collection
drives. Real lists, real trust, little bureaucracy. \\
5 & A Facebook blood group &
By far the most donors, and the least institutional weight. Useful for numbers,
weak as a story. \\
\bottomrule
\end{tabularx}
\end{center}

\note{Two different firsts. FADS is the best \emph{fit}; a university club is
the fastest \emph{yes}. Ask both --- a club that says yes this week is a
reference when you go to the committee next month, and neither call costs
anything.}

\subsection{None of this is a novel request}

The ANS publishes a page listing the associations it already collects blood
with: the FADS, the Scouts Musulmans Alg\'eriens, the Croissant-Rouge Alg\'erien,
Irchad wa Islah, the Rotary and Lions clubs, the UNJA and several others. Some
of those partnerships are dated on the page and go back to 2015.

That matters twice. It means an association hosting a collection is ordinary
rather than exceptional, so you are asking for something the organisation
already does. And it means the premise this product is built on --- that
associations are the right partner for reaching donors --- is the agency's own
practice, not an assumption you made.

% ===========================================================================
\section{What the ANS already publishes}

Before contacting anybody, three things are public and worth having in front of
you. The site is served over \texttt{http}; \texttt{https} currently fails on an
expired certificate, so a browser may refuse it.

\begin{center}
\begin{tabularx}{\textwidth}{@{}lX@{}}
\toprule
\textbf{Page} & \textbf{What is on it} \\
\midrule
\texttt{/centres-collectes-fixes/} &
A national directory: 165 transfusion structures across all 58 wilayas, each
with a telephone number. In Blida: CWTS CHU Blida 025 30 96 11, PTS EPH Meftah
025 33 60 04, BS EPH Boufarik 025 47 14 10. \\
\texttt{/programme-de-collectes-mobiles/} &
A mobile-collection programme published per wilaya. Dated 2025 at the time of
writing, so check whether the year has turned over before relying on it. \\
\texttt{/associations-donneur-sang/} &
The partnership page named above. Worth reading before any meeting: it tells
you what the agency already does and with whom. \\
\bottomrule
\end{tabularx}
\end{center}

\note{Agence Nationale du Sang --- (+213) 023 53 32 79 / 76 / 74, 15 route du
Kadous, Tixeraine, Bir Khadem, Alger. The number is printed on the back of
their mobile collection units.}

% ===========================================================================
\section{Where to go: a blood drive}

This is the single highest-yield place in the country for what you need, and
the reason is simple. \textbf{Everyone at a blood drive has already said yes.}
They are willing, they are present, and they are standing still.

You get the committee and the donors in the same room on the same afternoon.
Nowhere else does that.

\subsection{What to bring}

\begin{itemize}[leftmargin=1.4em]
  \item A phone with the app installed to the home screen, signed in.
  \item The invite code shown full screen --- \textbf{Committee $\rightarrow$
        Invite your donors $\rightarrow$ the QR button}. Hold it up, or hand
        the phone over.
  \item The same page printed on A4, taped to the table. The print button on
        that screen produces exactly this sheet and nothing else.
  \item Nothing to sign, no forms, no laptop.
\end{itemize}

\note{The code has no letter O, no digit 0, no I and no 1 --- it was designed to
be read aloud across a table when somebody's camera will not focus.}

% ===========================================================================
\section{What to say}

\subsection{The opening}

Do not begin with the application. Begin with the thing they already know is
broken.

\begin{say}
``When you need O negative on a Thursday night, what do you actually do?''
\end{say}

Let them answer. They will describe ringing round, or posting on Facebook, or a
WhatsApp group that half the members have muted. That answer is the product's
whole reason to exist, and it is better coming from them.

\subsection{The ask, which must be small}

The mistake is asking somebody to adopt a platform. Ask for an afternoon.

\begin{say}
``Can I put a code on your table for one afternoon? The people who scan it
become a list you can call --- and the app knows their blood type and the date
they are allowed to give again.''
\end{say}

They commit nothing. They install nothing. They keep the list.

\subsection{What is in it for them, said plainly}

\begin{say}
``Next time you need a specific type, you press one button. Only donors in your
wilaya, only compatible types, only the ones eligible today --- and never the
family who posted the request. Then you see who is coming.''
\end{say}

That last sentence is the one that lands. Nothing they use now tells them who is
coming.

% ===========================================================================
\section{The questions you will be asked}

Answer these plainly. The weak answers are the defensive ones.

\begin{description}[leftmargin=0em, style=unboxed]

\item[``Who are you?''] A student who built this, alone, and is looking for one
committee to try it. Do not inflate this. Nobody is offended by a student with
something that works; several people are offended by a student pretending to be
a company.

\item[``Is this official? Approved by the Agence Nationale du Sang?''] No, and
say so immediately. It is built to the law --- consent recorded per purpose,
every contact reveal logged, no payment to donors ever --- but nobody has
approved it. Then add the true thing: the ANS publishes a page of the
associations it collects blood with, and what is being proposed is that your
committee works the same way it already does, with a tool that remembers who
gave and when. Claiming an endorsement you do not have is the one mistake that
cannot be recovered from; pointing at their own published practice is not
claiming one.

\item[``What happens to our donors' phone numbers?''] They stay with the donor.
Search shows a masked number; opening one is a deliberate act, restricted to
members of an association in that wilaya, and every opening is written to a log
the donor can read and \emph{nobody} can delete. You can show them this on the
phone.

\item[``Are you taking our list?''] No, and this is worth being precise about.
There is no upload. Each donor signs up themselves and consents themselves; the
code only records that they came through your committee. If they walk away
tomorrow, the list is still theirs, because it was never yours to take.

\item[``We already have a Facebook group.''] So does everyone, and it works
until it does not. A group post reaches whoever is scrolling, cannot tell A+
from O$-$, does not know who gave blood six weeks ago and cannot be closed when
the need is met. This does all four.

\item[``Who pays for this?''] Nobody, today. Donor payment is excluded by
Algerian law and by the design --- there is no charge to a donor and there never
will be. Long term it is paid for by the bodies that coordinate donation, or by
public funding. Say that it is unsolved rather than inventing a business model
in front of somebody who will remember.

\item[``What if it stops working?''] Then they are exactly where they are now,
with a phone and a list. Nothing they do today is taken away.

\end{description}

% ===========================================================================
\section{On the day: the mechanics}

Ten minutes, and none of it needs a laptop.

\begin{center}
\begin{tabularx}{\textwidth}{@{}rX@{}}
\toprule
& \textbf{Step} \\
\midrule
1 & The committee's representative signs up in the app like anyone else. \\
2 & They apply as an association: \textbf{Committee $\rightarrow$ Apply}.
     Name, type, wilaya. \\
3 & You verify them. You are the platform administrator, and verification is
     restricted to that role on purpose --- a badge anybody could issue is not
     a badge. \\
4 & They open \textbf{Committee $\rightarrow$ Invite your donors} and create
     a link. Give it a name they will recognise later, like the date of the
     drive. \\
5 & Press the QR button. Hold up the phone, or print the sheet and tape it
     down. \\
6 & Each donor scans, signs up, chooses their blood type, and is on the list.
     The count on the committee's screen goes up as it happens. \\
\bottomrule
\end{tabularx}
\end{center}

\note{Creating a link is restricted to an administrator of that association,
because it publishes a standing invitation under the committee's name.
Volunteers can see everything and create nothing.}

% ===========================================================================
\section{What success looks like}

In order, and the second matters more than the first.

\begin{enumerate}[leftmargin=1.4em]
  \item \textbf{Twenty donors on one committee's list.} The count is on their
        screen; no spreadsheet needed.
  \item \textbf{One real request, answered.} A family posts, compatible donors
        are notified, somebody says they are coming, and the committee watches
        it happen. That is the entire product, and once it has happened once it
        has been proven.
  \item \textbf{A number you did not have before.} The committee screen reports
        how many requests reached a donor and how long the first one took.
        Today, on the demonstration data, it reads \emph{1 of 4, 38 hours}. In a
        real wilaya with real donors it should read minutes --- and if it does
        not, that is worth knowing more than any feature.
\end{enumerate}

% ===========================================================================
\section{What not to do}

\begin{itemize}[leftmargin=1.4em]
  \item \textbf{Do not promise a partnership you do not have.} Not with the
        Croissant-Rouge, not with the ANS, not with a hospital. One checkable
        exaggeration discredits every honest figure you have.
  \item \textbf{Do not ask them to install anything.} A link, scanned, opening
        in a browser. Anything heavier and half the room stops.
  \item \textbf{Do not demonstrate the sign-up with a real phone number that is
        not yours.} The demonstration build accepts one fixed verification
        code, which means anybody can claim any number. Use your own.
  \item \textbf{Do not collect names on paper ``to enter later''.} The entire
        argument is that nobody uploads anybody. Breaking that at the first
        drive would be worse than getting no donors at all.
\end{itemize}

\vfill
\note{Companion document: \emph{Qatra --- Technical and Strategic Dossier},
which carries the evidence, the regulatory position and the demonstration
script. Repository: \url{https://github.com/2him29/startup-idea} ---
demonstration at \url{https://2him29.github.io/startup-idea/}}

\end{document}
